\documentclass[10pt,a4paper]{article}
\usepackage[margin=2.2cm]{geometry}
\usepackage{fontspec}
\setmonofont{DejaVu Sans Mono}[Scale=0.84]
\newfontfamily\unifont{DejaVu Serif}[Scale=0.9]
\usepackage{amsmath,amssymb,mathtools}
\setcounter{MaxMatrixCols}{20}
\usepackage{longtable,booktabs,array}
\usepackage{fancyvrb}
\usepackage[dvipsnames]{xcolor}
\usepackage[colorlinks=true,linkcolor=NavyBlue,urlcolor=NavyBlue]{hyperref}
\usepackage{microtype}
\usepackage{needspace}
\usepackage{tocloft}\setlength{\cftbeforesecskip}{0.45em}
\setlength{\parindent}{0pt}\setlength{\parskip}{0.55em}
\setlength{\emergencystretch}{8em}
\renewcommand{\arraystretch}{1.15}
\allowdisplaybreaks
\begin{document}
\title{Gravity on the 7+1 spacetime: symbols and their definitions\\[0.4em]\large A dictionary of every symbol used for the curvature, the Einstein and Lovelock tensors, the generalized frame matrices, the canonical spin connection, the covariant Dirac equation of dirac16complex and the three cosmological epochs — sections 11–14 of the notebook 7+1.nb}
\author{Repository: github.com/once-ere/7plus1 — prepared for Nash\\ Notebook: 7+1.nb (Wolfram Language 15.0.1)\\ Document generated by tools/make\_docs.py on 2026-10-27}
\date{}
\maketitle
\begin{abstract}
For each symbol this dictionary gives the mathematical notation, the name used in the Mathematica notebook 7+1.nb (where there is one), the defining equation and its meaning. The symbols are grouped by topic in the order in which they are introduced; longer defining equations follow each table as displayed formulas. The last sections list symbols that look alike and the Mathematica names defined in sections 11–14 of 7+1.nb. Every equation needed to define a symbol is written in this document.
\end{abstract}
\tableofcontents
\Needspace{12\baselineskip}\section*{1. Conventions}\addcontentsline{toc}{section}{1. Conventions}
\begin{itemize}
\item Units: $\hbar=c=1$. A prime is $d/dx_0$ for functions of the time $x_0$ (and $d/dr$ for the functions $f(r)$, $h(r)$ of the static radial metric (2.4)).
\item Index ranges: Greek $\mu,\nu,\lambda,\rho,\sigma,\alpha,\beta,\gamma,\delta\in\lbrace0,\dots,7\rbrace$ are coordinate indices; Latin capitals $A,B,C,D\in\lbrace0,\dots,7\rbrace$ are frame indices; $i,j\in\lbrace1,2,3\rbrace$ label ordinary space and $k,l\in\lbrace4,\dots,7\rbrace$ the auxiliary space (for the static radial metric (2.4), $k\in\lbrace2,\dots,7\rbrace$). Repeated indices are summed.
\item The abstract indices of the xAct computation (section 4.3) are the letters $a$, $b$, $c$, $d$, $f$, $i$, $j$, $k$, $l$, $m$, $n$, $p$, $q$, $r$, $s$, $y$; they are unrelated to the index ranges above.
\item Frame indices are raised and lowered with $\eta$, coordinate indices with $g_{\mu\nu}$.
\item Matrices: $1_n$ is the $n\times n$ identity, $M^{\mathsf T}$ the transpose, $[X,Y]=XY-YX$, $\lbrace X,Y\rbrace=XY+YX$.
\item Mathematica lists are 1-based: $g_{\mu\nu}$ is \texttt{g[[μ+1, ν+1]]}, $\Gamma_A$ is \texttt{Γ17[A]} (an argument, not a list index). Greek letters in names are named characters: \texttt{\textbackslash{}[CapitalGamma]} is \texorpdfstring{\ensuremath{\Gamma}}{Γ}, \texttt{\textbackslash{}[Eta]} is \texorpdfstring{\ensuremath{\eta}}{η}, \texttt{\textbackslash{}[Alpha]} is \texorpdfstring{\ensuremath{\alpha}}{α}, \texttt{\textbackslash{}[Beta]} is \texorpdfstring{\ensuremath{\beta}}{β}, \texttt{\textbackslash{}[Delta]} is \texorpdfstring{\ensuremath{\delta}}{δ}, \texttt{\textbackslash{}[CapitalLambda]} is \texorpdfstring{\ensuremath{\Lambda}}{Λ}.
\item Christoffel symbols are written $\{{}^{\lambda}_{\mu\nu}\}$, never $\Gamma$: in this project $\Gamma$ always denotes a Dirac matrix.
\end{itemize}
\Needspace{12\baselineskip}\section*{2. Spacetime, metrics and coordinates}\addcontentsline{toc}{section}{2. Spacetime, metrics and coordinates}
{\small
\begin{longtable}{>{\raggedright\arraybackslash}p{\dimexpr 0.1400\linewidth-2\tabcolsep\relax}>{\raggedright\arraybackslash}p{\dimexpr 0.2400\linewidth-2\tabcolsep\relax}>{\raggedright\arraybackslash}p{\dimexpr 0.3500\linewidth-2\tabcolsep\relax}>{\raggedright\arraybackslash}p{\dimexpr 0.2700\linewidth-2\tabcolsep\relax}}
\toprule
\textbf{symbol} & \textbf{Mathematica} & \textbf{definition} & \textbf{meaning} \\
\midrule\endfirsthead
\toprule
\textbf{symbol} & \textbf{Mathematica} & \textbf{definition} & \textbf{meaning} \\
\midrule\endhead
\endfoot
\bottomrule\endlastfoot
$X$; $x_0,\dots,x_7$ & \texttt{X = \{x0, …, x7\}} & $X=(x_0,x_1,\dots,x_7)$ & $x_0$ time (the octonion unit direction), $x_1,x_2,x_3$ ordinary space, $x_4,\dots,x_7$ auxiliary space \\
$\eta_{AB}$ & \texttt{η17} & $\eta=\mathrm{diag}(+1,-1,-1,-1,-1,-1,-1,-1)$ & frame metric; signature $(+,-,\dots,-)$ \\
$g_{\mu\nu}$, $g^{\mu\nu}$ & \texttt{cv["g"]}, \texttt{cv["ginv"]} & $g^{\mu\lambda}g_{\lambda\nu}=\delta^\mu_\nu$ & metric and inverse metric \\
$|g|$ & — & $|\det g_{\mu\nu}|$ & volume factor $\sqrt{|g|}$ \\
$a$, $b$ & \texttt{a3[x0]}, \texttt{b4[x0]} & cosmological metric (2.1) & scale factors of ordinary and auxiliary space \\
$H_a$, $H_b$ & \texttt{Ha}, \texttt{Hb} & $H_a=a'/a$, $H_b=b'/b$ & Hubble rates \\
$H_a'$, $H_b'$ & \texttt{dHa}, \texttt{dHb} & derivatives of the Hubble rates & \texttt{toH} rewrites $a'$, $a''$, $b'$, $b''$ with them \\
$N$ & \texttt{nL[x0]} & $g_{00}=N^2$ & lapse function (variational check) \\
$n,a,b$ (inhomogeneous) & \texttt{nP}, \texttt{aP}, \texttt{bP} & polynomials (2.2) & inhomogeneous test metric \\
$F$ & \texttt{fPP} & arbitrary function of $(x_0-x_1,x_2,x_3)$ & profile of the plane wave (2.3) \\
$l_\mu$ & \texttt{lNull} & $l_\mu=(1,-1,0,\dots,0)$, $\eta^{\mu\nu}l_\mu l_\nu=0$ & null vector of the plane wave \\
$r$; $f$, $h$ & \texttt{x1}; \texttt{fR[x1]}, \texttt{hR[x1]} & static radial metric (2.4) & radial coordinate and two arbitrary positive functions \\
$C$ & \texttt{cConf[x0]} & $g=C(x_0)^2\eta$ & conformal factor \\
$H$, $q_H$ & \texttt{H}, \texttt{qH} & $a=e^{Hx_0}$, $b=e^{q_HHx_0}$ & constant Hubble rate and ratio of the exponential vacuum solutions (4.9); $H$ is also the inverse length of the field equations (4.7) \\*
\end{longtable}}
\textbf{(2.1) Cosmological metric:}
\begin{equation*}
ds^2=dx_0^2-a(x_0)^2\big(dx_1^2+dx_2^2+dx_3^2\big)-b(x_0)^2\big(dx_4^2+dx_5^2+dx_6^2+dx_7^2\big).
\end{equation*}
\textbf{(2.2) Inhomogeneous metric} (checked exactly at the point $x_0=1/2$, $x_4=1/3$, code \texttt{ptP}):
\begin{equation*}
ds^2=n^2dx_0^2-a^2\sum_{i=1}^3dx_i^2-b^2\sum_{k=4}^7dx_k^2,\quad n=1+\tfrac{x_0x_4}{5}+\tfrac{x_0^2}{11},\ \ a=1+\tfrac{x_0}{3}+\tfrac{x_4^2}{7},\ \ b=2-\tfrac{x_0x_4}{4}+\tfrac{x_0^2}{9}+\tfrac{x_4}{13}.
\end{equation*}
\textbf{(2.3) Plane gravitational wave (pp-wave):}
\begin{equation*}
ds^2=dx_0^2-dx_1^2-\cdots-dx_7^2+F(x_0-x_1,x_2,x_3)\,(dx_0-dx_1)^2 .
\end{equation*}
\textbf{(2.4) Static metric with a radial coordinate} $r=x_1>0$:
\begin{equation*}
ds^2=f(r)\,dx_0^2-h(r)\,dr^2-r^2\big(dx_2^2+\cdots+dx_7^2\big).
\end{equation*}
\Needspace{12\baselineskip}\section*{3. Curvature}\addcontentsline{toc}{section}{3. Curvature}
{\small
\begin{longtable}{>{\raggedright\arraybackslash}p{\dimexpr 0.1400\linewidth-2\tabcolsep\relax}>{\raggedright\arraybackslash}p{\dimexpr 0.2400\linewidth-2\tabcolsep\relax}>{\raggedright\arraybackslash}p{\dimexpr 0.3500\linewidth-2\tabcolsep\relax}>{\raggedright\arraybackslash}p{\dimexpr 0.2700\linewidth-2\tabcolsep\relax}}
\toprule
\textbf{symbol} & \textbf{Mathematica} & \textbf{definition} & \textbf{meaning} \\
\midrule\endfirsthead
\toprule
\textbf{symbol} & \textbf{Mathematica} & \textbf{definition} & \textbf{meaning} \\
\midrule\endhead
\endfoot
\bottomrule\endlastfoot
$\{{}^{\lambda}_{\mu\nu}\}$ & \texttt{cv["chr"]} & (3.1) & Christoffel symbols; \texttt{chr[[λ+1, μ+1, ν+1]]} \\
$R^\rho{}_{\sigma\mu\nu}$ & \texttt{cv["riem"]} & (3.2) & Riemann tensor (Landau–Lifshitz convention) \\
$R_{\rho\sigma\mu\nu}$ & \texttt{riemLow[cv]} & $g_{\rho\lambda}R^\lambda{}_{\sigma\mu\nu}$ & Riemann tensor, all indices down \\
$R^{\gamma\delta}{}_{\alpha\beta}$ & \texttt{cv["rud"]}, \texttt{rudAssoc} & $g^{\delta\sigma}R^\gamma{}_{\sigma\alpha\beta}$ & two indices up, two down; \texttt{rudAssoc} lists the non-zero entries with $\gamma<\delta$, $\alpha<\beta$ \\
$R_{\sigma\nu}$ & \texttt{cv["ric"]} & $R^\rho{}_{\sigma\rho\nu}$ & Ricci tensor \\
$R$ & \texttt{cv["R"]} & $g^{\sigma\nu}R_{\sigma\nu}$ & scalar curvature \\
$G^\mu{}_\nu$ & \texttt{cv["G"]} & $R^\mu{}_\nu-\tfrac12\delta^\mu_\nu R$ & Einstein tensor \\
$K_{AB}$ & \texttt{secK[cv, A, B]} & $R^{AB}{}_{AB}$ (no sum) & sectional curvature of the cosmological metric \\
$\nabla_\mu T^\mu{}_\nu$ & \texttt{divergenceMixed} & (3.3) & divergence of a mixed tensor \\
— & \texttt{curvature[g, simp]} & — & the curvature engine: returns all of the above for a metric matrix \\*
\end{longtable}}
\begin{align*}
\{{}^{\lambda}_{\mu\nu}\}&=\tfrac12\,g^{\lambda\rho}\big(\partial_\mu g_{\rho\nu}+\partial_\nu g_{\rho\mu}-\partial_\rho g_{\mu\nu}\big),&&(3.1) \\
R^{\rho}{}_{\sigma\mu\nu}&=\partial_\mu\{{}^{\rho}_{\sigma\nu}\}-\partial_\nu\{{}^{\rho}_{\sigma\mu}\}+\{{}^{\rho}_{\mu\lambda}\}\{{}^{\lambda}_{\sigma\nu}\}-\{{}^{\rho}_{\nu\lambda}\}\{{}^{\lambda}_{\sigma\mu}\},&&(3.2) \\
\nabla_\mu T^\mu{}_\nu&=\partial_\mu T^\mu{}_\nu+\{{}^{\mu}_{\mu\lambda}\}T^\lambda{}_\nu-\{{}^{\lambda}_{\mu\nu}\}T^\mu{}_\lambda .&&(3.3)
\end{align*}
\Needspace{12\baselineskip}\section*{4. Lovelock tensors and field equations}\addcontentsline{toc}{section}{4. Lovelock tensors and field equations}
\Needspace{8\baselineskip}\subsection*{4.1 Generalized Kronecker delta and Lovelock tensors}\addcontentsline{toc}{subsection}{4.1 Generalized Kronecker delta and Lovelock tensors}
{\small
\begin{longtable}{>{\raggedright\arraybackslash}p{\dimexpr 0.1400\linewidth-2\tabcolsep\relax}>{\raggedright\arraybackslash}p{\dimexpr 0.2400\linewidth-2\tabcolsep\relax}>{\raggedright\arraybackslash}p{\dimexpr 0.3500\linewidth-2\tabcolsep\relax}>{\raggedright\arraybackslash}p{\dimexpr 0.2700\linewidth-2\tabcolsep\relax}}
\toprule
\textbf{symbol} & \textbf{Mathematica} & \textbf{definition} & \textbf{meaning} \\
\midrule\endfirsthead
\toprule
\textbf{symbol} & \textbf{Mathematica} & \textbf{definition} & \textbf{meaning} \\
\midrule\endhead
\endfoot
\bottomrule\endlastfoot
$\delta^{\mu_1\cdots\mu_n}_{\nu_1\cdots\nu_n}$ & \texttt{gKD[up, dn]} & (4.1) & generalized Kronecker delta (explicit index values) \\
$\delta^{\text{upper}}_{\text{lower}}$ & \texttt{kδ[lower, upper]} & (4.1) with xAct \texttt{delta} & the same with abstract indices (xAct) \\
$E^{(k)\mu}{}_\nu$ & \texttt{lovelockE[ra, k]}; \texttt{Lovelock1/2/3} & (4.2) & Lovelock tensor of order $k$; \texttt{Lovelock}$k$ $=E^{(k)b}{}_a$ \\
$L_k$ & \texttt{lovelockL[ra, k]} & (4.3) & Lovelock density of order $k$ \\
$E^{(0)}$ & \texttt{lovelockE[ra, 0]} & $-\tfrac12\delta^\mu_\nu$ & cosmological term \\
$L_2$ & \texttt{lanczosDiag[cv][[1]]} & $R^2-4R_{\mu\nu}R^{\mu\nu}+R_{\mu\nu\rho\sigma}R^{\mu\nu\rho\sigma}$ & Gauss–Bonnet density \\
$H_{\mu\nu}$ & \texttt{lanczosDiag[cv][[2]]} & (4.4) & Lanczos tensor ($=E^{(2)}$) \\
$\mathcal U$, $\mathcal V$ & \texttt{pairSets}, \texttt{orderedSeqs} & sets of $k$ disjoint pairs; ordered sequences of pairs & index pairs of the pair-set algorithm (4.5) \\
sgn & \texttt{Signature[up] Signature[dn]} & sign of the permutation upper → lower &  \\
$E^{(k)}$ (cosmological) & \texttt{ECos[[k+1]]}, \texttt{LCos[[k+1]]} & closed forms in $H_a,H_b,H_a',H_b'$ & Lovelock tensors and densities of (2.1) \\*
\end{longtable}}
\begin{align*}
\delta^{\mu_1\cdots\mu_n}_{\nu_1\cdots\nu_n}&=\det\big[\delta^{\mu_i}_{\nu_j}\big]_{i,j=1}^{n},&&(4.1) \\
E^{(k)\mu}{}_{\nu}&=-\frac{1}{2^{k+1}}\,\delta^{\mu\alpha_1\beta_1\cdots\alpha_k\beta_k}_{\nu\gamma_1\delta_1\cdots\gamma_k\delta_k}\,R^{\gamma_1\delta_1}{}_{\alpha_1\beta_1}\cdots R^{\gamma_k\delta_k}{}_{\alpha_k\beta_k},&&(4.2) \\
L_k&=\frac{1}{2^{k}}\,\delta^{\alpha_1\beta_1\cdots\alpha_k\beta_k}_{\gamma_1\delta_1\cdots\gamma_k\delta_k}\,R^{\gamma_1\delta_1}{}_{\alpha_1\beta_1}\cdots R^{\gamma_k\delta_k}{}_{\alpha_k\beta_k},&&(4.3) \\
H_{\mu\nu}&=2\big(RR_{\mu\nu}-2R_{\mu\alpha}R^\alpha{}_\nu-2R^{\alpha\beta}R_{\mu\alpha\nu\beta}+R_\mu{}^{\alpha\beta\gamma}R_{\nu\alpha\beta\gamma}\big)-\tfrac12g_{\mu\nu}L_2,&&(4.4) \\
E^{(k)\mu}{}_\nu&=-2^{k-1}k!\sum_{\mathcal U}\sum_{\mathcal V}\mathrm{sgn}(\mu\,\mathcal U\to\nu\,\mathcal V)\prod_{i=1}^kR^{\mathcal V_i}{}_{\mathcal U_i}.&&(4.5)
\end{align*}
\Needspace{8\baselineskip}\subsection*{4.2 The xAct objects (the method of the reference notebook Einstein-Lovelock-7+1-Nash.nb)}\addcontentsline{toc}{subsection}{4.2 The xAct objects (the method of the reference notebook Einstein-Lovelock-7+1-Nash.nb)}
{\small
\begin{longtable}{>{\raggedright\arraybackslash}p{\dimexpr 0.1300\linewidth-2\tabcolsep\relax}>{\raggedright\arraybackslash}p{\dimexpr 0.2200\linewidth-2\tabcolsep\relax}>{\raggedright\arraybackslash}p{\dimexpr 0.4000\linewidth-2\tabcolsep\relax}>{\raggedright\arraybackslash}p{\dimexpr 0.2500\linewidth-2\tabcolsep\relax}}
\toprule
\textbf{symbol} & \textbf{Mathematica} & \textbf{definition} & \textbf{meaning} \\
\midrule\endfirsthead
\toprule
\textbf{symbol} & \textbf{Mathematica} & \textbf{definition} & \textbf{meaning} \\
\midrule\endhead
\endfoot
\bottomrule\endlastfoot
$M_8$ & \texttt{M8} & \texttt{DefManifold[M8, 8, \{a, b, c, d, f, i, j, k, l, m, n, p, q, r, s, y\}]} & the 8-dimensional manifold with its abstract indices \\
$g$ & \texttt{g71} & \texttt{DefMetric[\{1, 7, 0\}, g71[-a, -b], CD, …]} & abstract metric; $\{1,7,0\}$ = one plus, seven minus signs \\
$\nabla$ & \texttt{CD} & Levi-Civita derivative of \texttt{g71} & after \texttt{CD = LC[g71]}: the connection of a component metric \\
chart & \texttt{chartM8} & \texttt{DefChart[chartM8, M8, \{0, …, 7\}, \{x0[], …, x7[]\}, …]} & coordinate chart; \texttt{x0[]} … \texttt{x7[]} coordinate scalars \\
$\delta_a{}^b$ & \texttt{delta[-a, b]} & Kronecker delta of xAct &  \\
$R_{ab}{}^{cd}$ etc. & \texttt{RiemannCD}, \texttt{RicciCD}, \texttt{RicciScalarCD}, \texttt{EinsteinCD} & curvature of \texttt{CD} (abstract, or components after the assignments) &  \\
— & \texttt{MyArrayComponents} & \texttt{MyArrayComponents[expr] := expr // ToBasis[chartM8] // ComponentArray // ToValues // ToValues // Simplify} & component array of an expression \\
— & \texttt{toXc}, \texttt{fromXc} & $x_\mu\to$ \texttt{xμ[]} and back & between this notebook’s coordinates and the chart scalars \\
— & \texttt{xActLovelock[gm]} & \texttt{g71 = CTensor[gm, \{-chartM8, -chartM8\}]}; \texttt{MetricCompute}; \texttt{CD = LC[g71]}; \texttt{RiemannCD = Riemann[CD]}, …; \texttt{MyArrayComponents} of \texttt{Lovelock1/2/3}, \texttt{EinsteinCD[-a, b]}, transposed & returns $E^{(1)},E^{(2)},E^{(3)},G$ as matrices (row $\mu$, column $\nu$) \\
— & \texttt{xlCos}, \texttt{xlPP}, \texttt{xlExp} & \texttt{xActLovelock} of (2.1), (2.3) and of $a=e^{Hx_0}$, $b=e^{q_HHx_0}$ &  \\*
\end{longtable}}
\Needspace{8\baselineskip}\subsection*{4.3 Couplings and field equations}\addcontentsline{toc}{subsection}{4.3 Couplings and field equations}
{\small
\begin{longtable}{>{\raggedright\arraybackslash}p{\dimexpr 0.1400\linewidth-2\tabcolsep\relax}>{\raggedright\arraybackslash}p{\dimexpr 0.2400\linewidth-2\tabcolsep\relax}>{\raggedright\arraybackslash}p{\dimexpr 0.3500\linewidth-2\tabcolsep\relax}>{\raggedright\arraybackslash}p{\dimexpr 0.2700\linewidth-2\tabcolsep\relax}}
\toprule
\textbf{symbol} & \textbf{Mathematica} & \textbf{definition} & \textbf{meaning} \\
\midrule\endfirsthead
\toprule
\textbf{symbol} & \textbf{Mathematica} & \textbf{definition} & \textbf{meaning} \\
\midrule\endhead
\endfoot
\bottomrule\endlastfoot
$\kappa$, $G_8$ & \texttt{kap} & $\kappa=8\pi G_8$ & gravitational coupling; Newton constant in 8 dimensions \\
$\alpha_k$ & \texttt{alf[k]} & (4.6) & Lovelock couplings, $\alpha_1=1$, $\alpha_2$ (length²), $\alpha_3$ (length\texorpdfstring{\ensuremath{{}^{4}}}{⁴}) \\
$\bar\Lambda$ & \texttt{lamb} & $\alpha_0=2\bar\Lambda$ & cosmological constant (length$^{-2}$); in the field equations of the three epochs (section 7) it is written $\Lambda$ \\
$w_1,w_2,w_3$ & \texttt{w1}, \texttt{w2}, \texttt{w3} & (4.7) & pure-number couplings of the reference notebook \\
$\Lambda$ & \texttt{\textbackslash{}[CapitalLambda]} & (4.7) & pure-number cosmological constant of the reference notebook \\
$H$ & \texttt{H} & (4.7) & fundamental inverse length \\
$T^\mu{}_\nu$ & — & $\delta S_{\rm matter}=\tfrac12\int\sqrt{|g|}T_{\mu\nu}\delta g^{\mu\nu}$ & energy–momentum tensor \\
$\varepsilon$, $p_a$, $p_b$ & — & $T=\mathrm{diag}(\varepsilon,-p_a,-p_a,-p_a,-p_b,\dots,-p_b)$ & energy density and pressures \\*
\end{longtable}}
\begin{align*}
S&=-\frac{1}{2\kappa}\int d^8X\sqrt{|g|}\sum_{k=0}^3\alpha_kL_k+S_{\rm matter},\qquad \sum_{k=0}^{3}\alpha_kE^{(k)\mu}{}_\nu=\kappa T^\mu{}_\nu,&&(4.6) \\
0&=-\Lambda\,\delta^\mu_\nu+H^{-2}w_1\,E^{(1)\mu}{}_\nu+H^{-4}w_2\,E^{(2)\mu}{}_\nu+H^{-6}w_3\,E^{(3)\mu}{}_\nu,&&(4.7) \\
\alpha_2&=\frac{w_2}{w_1H^2},\quad \alpha_3=\frac{w_3}{w_1H^4},\quad \bar\Lambda=\frac{\Lambda H^2}{w_1}\quad(T=0),&&(4.8) \\
a&=e^{Hx_0},\quad b=e^{q_HHx_0}:\quad H_a=H,\ H_b=q_HH,\ H_a'=H_b'=0;&&(4.9) \\
&\qquad q_H=-1:\ \ w_2=\tfrac{\Lambda}{4}-\tfrac{w_1}{2},\ \ w_3=\tfrac{w_1}{24}-\tfrac{\Lambda}{36}.
\end{align*}
\Needspace{12\baselineskip}\section*{5. Frames, spin connection, spinor curvature}\addcontentsline{toc}{section}{5. Frames, spin connection, spinor curvature}
{\small
\begin{longtable}{>{\raggedright\arraybackslash}p{\dimexpr 0.1400\linewidth-2\tabcolsep\relax}>{\raggedright\arraybackslash}p{\dimexpr 0.2400\linewidth-2\tabcolsep\relax}>{\raggedright\arraybackslash}p{\dimexpr 0.3500\linewidth-2\tabcolsep\relax}>{\raggedright\arraybackslash}p{\dimexpr 0.2700\linewidth-2\tabcolsep\relax}}
\toprule
\textbf{symbol} & \textbf{Mathematica} & \textbf{definition} & \textbf{meaning} \\
\midrule\endfirsthead
\toprule
\textbf{symbol} & \textbf{Mathematica} & \textbf{definition} & \textbf{meaning} \\
\midrule\endhead
\endfoot
\bottomrule\endlastfoot
$e^A{}_\mu$ & \texttt{fd["e"]} & $g_{\mu\nu}=\eta_{AB}e^A{}_\mu e^B{}_\nu$ & frame field (vielbein); \texttt{e[[A+1, μ+1]]} \\
$e_A{}^\mu$ & \texttt{fd["ei"]} & $e_A{}^\mu e^A{}_\nu=\delta^\mu_\nu$ & inverse frame; \texttt{ei[[μ+1, A+1]]} \\
$\Gamma_A$ & \texttt{Γ17[A]} & $\Gamma_0=\left(\begin{smallmatrix}0&1_8\\1_8&0\end{smallmatrix}\right)$, $\Gamma_i=\left(\begin{smallmatrix}\tau^O_i&0\\0&-\tau^O_i\end{smallmatrix}\right)$ & real 16×16 Dirac matrices, $\lbrace\Gamma_A,\Gamma_B\rbrace=2\eta_{AB}1_{16}$ \\
$\Gamma^A$ & \texttt{Γ17up[A]} & $\eta^{AA}\Gamma_A$ (no sum) &  \\
$\Gamma_\mu(X)$ & \texttt{fd["Gam"]} & $e^A{}_\mu\Gamma_A$ & generalized frame matrices \\
$\Gamma^\mu(X)$ & \texttt{fd["GamUp"]} & $e_A{}^\mu\Gamma^A$ &  \\
$\omega_\mu{}^A{}_B$ & \texttt{fd["om"]} & (5.1) & canonical spin connection \\
$\omega_{\mu AB}$ & \texttt{fd["omL"]} & $\eta_{AC}\omega_\mu{}^C{}_B$ & antisymmetric in $A,B$ \\
$\Omega_\mu$ & \texttt{fd["Om"]} & $\tfrac14\omega_{\mu AB}\Gamma^A\Gamma^B$ & spinor connection (real 16×16) \\
$\mathcal F_{\mu\nu}$ & \texttt{fd["F"]} & (5.2) & spinor curvature \\
$R_{AB\mu\nu}$ & \texttt{rABmn} & $\eta_{AC}e^C{}_\rho R^\rho{}_{\sigma\mu\nu}e_B{}^\sigma$ & Riemann tensor with two frame indices \\
$C^A{}_{BC}$ & \texttt{structureC} & (5.3) & structure coefficients of the frame \\
$\omega_{C,AB}$ & \texttt{omegaFrame} & $e_C{}^\mu\omega_{\mu AB}$ & spin connection with a frame form index \\
$1_{16}$, $0_{16}$ & \texttt{ID16}, \texttt{Zero16} &  & identity and zero 16×16 matrices \\
— & \texttt{frameData}, \texttt{frameChecks}, \texttt{cartanEq}, \texttt{structureFormula} &  & compute and check the frame data of a metric \\*
\end{longtable}}
\begin{align*}
\omega_\mu{}^A{}_B&=e^A{}_\nu\big(\partial_\mu e_B{}^\nu+\{{}^{\nu}_{\mu\lambda}\}e_B{}^\lambda\big),&&(5.1) \\
\mathcal F_{\mu\nu}&=\partial_\mu\Omega_\nu-\partial_\nu\Omega_\mu+[\Omega_\mu,\Omega_\nu]=\tfrac14R_{AB\mu\nu}\Gamma^A\Gamma^B,&&(5.2) \\
C^A{}_{BC}&=e_B{}^\mu e_C{}^\nu\big(\partial_\mu e^A{}_\nu-\partial_\nu e^A{}_\mu\big),\qquad \omega_{C,AB}=-\tfrac12\big(C_{ACB}+C_{BAC}-C_{CBA}\big).&&(5.3)
\end{align*}
The total covariant derivative and the vielbein postulate:
\begin{equation*}
\nabla_\mu\Gamma_\nu=\partial_\mu\Gamma_\nu-\{{}^{\lambda}_{\mu\nu}\}\Gamma_\lambda+[\Omega_\mu,\Gamma_\nu]=0,\qquad \partial_\mu e^A{}_\nu-\{{}^{\lambda}_{\mu\nu}\}e^A{}_\lambda+\omega_\mu{}^A{}_Be^B{}_\nu=0 .
\end{equation*}
\Needspace{12\baselineskip}\section*{6. The covariant Dirac equation of dirac16complex}\addcontentsline{toc}{section}{6. The covariant Dirac equation of dirac16complex}
{\small
\begin{longtable}{>{\raggedright\arraybackslash}p{\dimexpr 0.1400\linewidth-2\tabcolsep\relax}>{\raggedright\arraybackslash}p{\dimexpr 0.2400\linewidth-2\tabcolsep\relax}>{\raggedright\arraybackslash}p{\dimexpr 0.3500\linewidth-2\tabcolsep\relax}>{\raggedright\arraybackslash}p{\dimexpr 0.2700\linewidth-2\tabcolsep\relax}}
\toprule
\textbf{symbol} & \textbf{Mathematica} & \textbf{definition} & \textbf{meaning} \\
\midrule\endfirsthead
\toprule
\textbf{symbol} & \textbf{Mathematica} & \textbf{definition} & \textbf{meaning} \\
\midrule\endhead
\endfoot
\bottomrule\endlastfoot
$\psi$ & \texttt{psiF} & 16 functions of $(x_0,x_1,x_5)$ in the checks & the field dirac16complex \\
$\bar\psi$ & \texttt{psibF} & $\psi^\dagger\Gamma_0$ (independent functions in the checks) & Dirac adjoint \\
$D_\mu$ & \texttt{DpsiCos}, \texttt{DpsibCos} & (6.1) & spinor covariant derivative \\
$m$ & \texttt{mD} &  & mass of dirac16complex \\
$J^\mu$ & — & $\bar\psi\Gamma^\mu\psi$ & conserved current \\
$\chi$ & \texttt{chiF} & $\chi=(a^3b^4)^{1/2}\psi$ & rescaled field (cosmological metric) \\
$c(x_0)$ & \texttt{cModes} & $\chi=e^{i(\vec p\cdot\vec x+\vec q\cdot\vec y)}c(x_0)$ & mode function (16 components) \\
$\vec p$, $\vec q$ & \texttt{pmom[i]}, \texttt{qmom[k]} & $\vec p=(p_1,p_2,p_3)$, $\vec q=(q_4,\dots,q_7)$ & comoving momenta along ordinary and auxiliary space \\
$\alpha_j$, $\beta$ & \texttt{αD[j]}, \texttt{βD} & $\alpha_j=\Gamma^0\Gamma^j$, $\beta=\Gamma_0$ & Hamiltonian Dirac matrices \\
$h(x_0)$ & \texttt{hCos} & (6.2) & mode Hamiltonian \\
$A_\mu$, $q$, $\alpha$ & \texttt{aGauge[m]}, \texttt{qG}, \texttt{alGauge} & (6.3) & abelian gauge field, charge, gauge function \\
$\chi$ (radial) & \texttt{chiR} & $\psi=f^{-1/4}r^{-3}\chi$ & rescaled field (static radial metric) \\*
\end{longtable}}
\begin{align*}
D_\mu\psi&=\partial_\mu\psi+\Omega_\mu\psi,\qquad D_\mu\bar\psi=\partial_\mu\bar\psi-\bar\psi\,\Omega_\mu,\qquad (i\Gamma^\mu D_\mu-m)\psi=0,&&(6.1) \\
h(x_0)&=\sum_{i=1}^{3}\alpha_i\frac{p_i}{a}+\sum_{k=4}^{7}\alpha_k\frac{q_k}{b}+\beta m,\qquad i\frac{dc}{dx_0}=h\,c,&&(6.2) \\
D_\mu\psi&=\partial_\mu\psi+\Omega_\mu\psi-iqA_\mu\psi,\qquad \psi\to e^{iq\alpha}\psi,\ \ A_\mu\to A_\mu+\partial_\mu\alpha .&&(6.3)
\end{align*}
\Needspace{12\baselineskip}\section*{7. The three epochs}\addcontentsline{toc}{section}{7. The three epochs}
{\small
\begin{longtable}{>{\raggedright\arraybackslash}p{\dimexpr 0.1400\linewidth-2\tabcolsep\relax}>{\raggedright\arraybackslash}p{\dimexpr 0.2400\linewidth-2\tabcolsep\relax}>{\raggedright\arraybackslash}p{\dimexpr 0.3500\linewidth-2\tabcolsep\relax}>{\raggedright\arraybackslash}p{\dimexpr 0.2700\linewidth-2\tabcolsep\relax}}
\toprule
\textbf{symbol} & \textbf{Mathematica} & \textbf{definition} & \textbf{meaning} \\
\midrule\endfirsthead
\toprule
\textbf{symbol} & \textbf{Mathematica} & \textbf{definition} & \textbf{meaning} \\
\midrule\endhead
\endfoot
\bottomrule\endlastfoot
$p$, $q$ & \texttt{pK}, \texttt{qK} & $a=x_0^p$, $b=x_0^q$ & Kasner exponents \\
$p_E$, $q_E$ & \texttt{pKE}, \texttt{qKE} & $p=(1+2\sqrt2)/7$, $q=(2-3\sqrt2)/14$ & Einstein–Kasner branch (expanding ordinary space, contracting auxiliary space) \\
$K$, $\lambda$ & — & $K=|\vec p|\lambda^{1-p}$, $\lambda=(|\vec q|/|\vec p|)^{1/(q-p)}$ & scaling variable of massless Kasner modes \\
$\mathcal A$, $\mathcal B$, $\mathcal Z$ & \texttt{Aop}, \texttt{Bop2}, \texttt{Zc} & $\mathcal A=\vec n_a\cdot\vec\alpha$, $\mathcal B=\vec n_b\cdot\vec\alpha$, $\mathcal Z=\mathcal A\mathcal B\beta$ & $\mathcal Z^2=-1_{16}$: conserved $\pm i$ \\
$\sigma_{1,2,3}$ & \texttt{pauli} & Pauli matrices & 2×2 reduction of the mode equation \\
$|\alpha|^2$, $|\beta|^2$ & \texttt{bogoliubov[h, t0, t1, n]} & returns $\lbrace|\beta|^2,|\alpha|^2\rbrace$ & Bogoliubov coefficients: pair-creation probability \\
$b_0$, $V_\perp$, $V_c$ & — & $V_\perp=b_0^4V_c$ & static auxiliary scale, its proper and coordinate volume \\
$G_4$ & — & $G_4=G_8/V_\perp$ & four-dimensional Newton constant \\
$M$ & — & $\sqrt{m^2+|\vec q|^2/b_0^2}$ & Kaluza–Klein mass \\
$H$ (de Sitter) & \texttt{hdS} & $a=e^{Hx_0}$ & de Sitter expansion rate \\
$T_{\rm GH}$ & — & $H/(2\pi)$ & Gibbons–Hawking temperature \\*
\end{longtable}}
\begin{equation*}
|\beta|^2=\frac{1}{e^{2\pi M/H}+1}\quad\text{(de Sitter)},\qquad |\beta|^2=e^{-\pi\kappa_{\rm LZ}^2}\quad\text{(Landau–Zener benchmark, } i\dot c=(\kappa_{\rm LZ}\sigma_1+t\sigma_3)c\text{)}.
\end{equation*}
\Needspace{12\baselineskip}\section*{8. Symbols that look alike}\addcontentsline{toc}{section}{8. Symbols that look alike}
{\small
\begin{longtable}{>{\raggedright\arraybackslash}p{\dimexpr 0.4200\linewidth-2\tabcolsep\relax}>{\raggedright\arraybackslash}p{\dimexpr 0.5800\linewidth-2\tabcolsep\relax}}
\toprule
\textbf{symbol} & \textbf{not to be confused with} \\
\midrule\endfirsthead
\toprule
\textbf{symbol} & \textbf{not to be confused with} \\
\midrule\endhead
\endfoot
\bottomrule\endlastfoot
$\Gamma_A$, $\Gamma_\mu$ (Dirac matrices) & $\{{}^{\lambda}_{\mu\nu}\}$ (Christoffel symbols; never written $\Gamma$ here) \\
$H$ (inverse length of (4.7); de Sitter rate) & $H_a$, $H_b$ (Hubble rates), $H_{\mu\nu}$ (Lanczos tensor) \\
$\Lambda$ (pure number, (4.7)) & $\bar\Lambda$ (cosmological constant with dimension length$^{-2}$) \\
$\alpha_k$ (Lovelock couplings) & $\alpha_j$ (Dirac matrices $\Gamma^0\Gamma^j$), $\alpha$ (gauge function), $|\alpha|^2$ (Bogoliubov) \\
$\beta$ (Dirac matrix $\Gamma_0$) & $|\beta|^2$ (Bogoliubov coefficient) \\
$q$ (charge; Kasner exponent of $b$; xAct index) & $q_H$ (ratio of Hubble rates, code \texttt{qH}), $\vec q$ (auxiliary momentum) \\
$w_1,w_2,w_3$ (Lovelock couplings) & $W^1,W^2,W^3$ (electroweak gauge fields, code \texttt{wW1}, \texttt{wW2}, \texttt{wW3}, in sections 17–19 of 7+1.nb) \\
$K$ (Kasner scaling variable) & $K_{AB}$ (sectional curvature) \\
$\kappa=8\pi G_8$ & $\kappa_{\rm LZ}$ (Landau–Zener coupling) \\
\texttt{Lovelock1} $=E^{(1)b}{}_a$ (xAct) & \texttt{lovelockE[ra, 1]} $=E^{(1)\mu}{}_\nu$ (pair-set sum): the same tensor, row/column order transposed by \texttt{xActLovelock} \\*
\end{longtable}}
\end{document}
